%=====================================================================
\chapter{Capstone: One Problem, Four Domains}\label{ch:capstone}
%=====================================================================
\locator{6}{Part VI --- Professional Practice}

\begin{outcomes}
By the end of this chapter you will be able to:
\begin{tight}
  \item \textbf{Execute} a complete data science project from framing to monitoring, unaided.
  \item \textbf{Transfer} one technical approach across four application domains, adapting it to each.
  \item \textbf{Justify} every design decision by reference to the constraint that forced it.
  \item \textbf{Produce} the four artefacts a professional project delivers.
  \item \textbf{Assess} your own work against the rubric used to mark it.
\end{tight}
\end{outcomes}

\begin{prereq}
All thirty-one preceding chapters. This one introduces no new technique --- it asks
whether you can choose among the ones you have.
\end{prereq}

\section{The Shared Problem}

Every Four Lenses panel in this book has hinted at something worth making explicit
now: the four domains have been solving \emph{the same problem} in different
clothing.

\dsfig{%
\begin{tikzpicture}[font=\scriptsize,
  core/.style={rounded corners=5pt, draw=primary, line width=1.4pt, fill=primary!8,
               text=primary, text width=4.6cm, align=center, font=\scriptsize\bfseries,
               inner sep=7pt},
  d/.style={rounded corners=3pt, draw=#1, fill=#1!8, text=#1!50!black,
            text width=3.3cm, align=center, font=\tiny, inner sep=5pt, minimum height=2.0cm},
  arr/.style={-{Stealth[length=2.2mm]}, primary!45, line width=0.9pt}]

\node[core] (c) at (5.75,0.4) {%
  THE COMMON PROBLEM\\[3pt]
  \normalfont\tiny Detect, early and from behavioural traces, which entities are
  heading for a bad outcome --- while there is still time to act, and with only as
  many alerts as someone can act on.};

\node[d=lensLA]  (la) at (0,-3.0)     {%
  {\color{lensLA}\bfseries\faIcon{graduation-cap}~Learning Analytics}\\[3pt]
  \textbf{Entity:} student\\
  \textbf{Bad outcome:} failure\\
  \textbf{Trace:} clickstream\\
  \textbf{Budget:} adviser hours};

\node[d=lensPM]  (pm) at (3.85,-3.0)  {%
  {\color{lensPM}\bfseries\faIcon{tasks}~Project Management}\\[3pt]
  \textbf{Entity:} sprint\\
  \textbf{Bad outcome:} overrun\\
  \textbf{Trace:} ticket activity\\
  \textbf{Budget:} attention};

\node[d=lensIOT] (io) at (7.7,-3.0)   {%
  {\color{lensIOT}\bfseries\faIcon{microchip}~Internet of Things}\\[3pt]
  \textbf{Entity:} motor\\
  \textbf{Bad outcome:} failure\\
  \textbf{Trace:} telemetry\\
  \textbf{Budget:} technicians};

\node[d=lensSEC] (se) at (11.55,-3.0) {%
  {\color{lensSEC}\bfseries\faIcon{shield-alt}~Cybersecurity}\\[3pt]
  \textbf{Entity:} account\\
  \textbf{Bad outcome:} compromise\\
  \textbf{Trace:} auth logs\\
  \textbf{Budget:} analysts};

\foreach \n in {la,pm,io,se}{\draw[arr] (c) -- (\n);}

\node[anchor=north west, align=left, text width=13.8cm, font=\scriptsize, text=accent!75!black]
     at (-2.0,-4.6)
     {\textbf{The differences that matter are not technical.} All four are rare-outcome
      detection on behavioural time series with a hard capacity constraint. What differs
      is the \emph{base rate} (from 8\% to 0.01\%), the \emph{ground-truth lag} (a term,
      a fortnight, a quarter, never), whether the entity \emph{knows} it is being modelled
      and adapts, and what happens to a person when you are wrong. Those four differences
      determine your design.};
\end{tikzpicture}}%
{The common structure beneath the four domains. Recognising it is the point of the whole
handout: the technique transfers, and the constraints do not.}{capstone-common}

\rowcolors{2}{white}{lightbg}
\begin{longtable}{@{}L{2.9cm}L{2.5cm}L{2.5cm}L{2.5cm}L{2.6cm}@{}}
\toprule
\rowcolor{primary!15}
& \textbf{Learning} & \textbf{Projects} & \textbf{IoT} & \textbf{Security} \\
\midrule
\endhead
Base rate & $\sim$8\% & $\sim$50\% & $\sim$0.2\% & $\sim$0.001\% \\
Rows available & thousands & \textbf{dozens} & millions & billions \\
Labels? & yes, lagged & yes, few & \textbf{rarely} & few, biased \\
Ground-truth lag & one term & one sprint & months & often never \\
Adversarial? & mildly & no & no & \textbf{yes} \\
Cost of a miss & a lost student & a late release & a line stoppage & a breach \\
Cost of a false alarm & a wasted call & lost credibility & a wasted inspection & analyst time \\
Fairness stakes & \textbf{very high} & high & low & moderate \\
Right primary metric & recall @ capacity & $F_1$ / interval & recall + lead time & precision @ capacity \\
Right model class & interpretable & \textbf{very simple} & boosting / anomaly & ensemble + anomaly \\
Right validation & walk-forward, by cohort & chronological & by device \& time & strictly chronological \\
\bottomrule
\end{longtable}

\begin{conceptbox}[Read That Table Downwards, Then Across]
Downwards, it is four independent project briefs. Across, it is the answer to
``what actually changes between domains?'' --- and the answer is almost never the
algorithm. It is the base rate, the labels, the lag, the adversary and the stakes.
A graduate who can reason across that table is employable in all four fields.
\end{conceptbox}

\section{The Capstone Brief}

\begin{definitionbox}[Your Task]
Choose \textbf{one} domain from \dsref{capstone-common}. Execute a complete project
on it, from framing to monitoring plan. Then write one page explaining how your
design would have to change if you moved it to a \emph{different} one of the four
--- naming the specific constraint that forces each change.
\end{definitionbox}

\dsfig{%
\begin{tikzpicture}[font=\scriptsize,
  st/.style={rounded corners=3pt, draw=#1, fill=#1!8, text=#1!50!black,
             text width=2.45cm, align=center, font=\tiny\bfseries, minimum height=0.85cm},
  arr/.style={-{Stealth[length=2mm]}, primary!55, line width=0.95pt},
  ch/.style={font=\tiny\itshape, text=gray!55!black, align=center, text width=2.6cm}]

\node[st=primary]   (a) at (0,0)      {1. FRAME\\{\tiny 6 questions}};
\node[st=secondary] (b) at (2.85,0)   {2. AUDIT\\{\tiny 6 dimensions}};
\node[st=tealhl]    (c) at (5.7,0)    {3. EXPLORE};
\node[st=goldhl]    (d) at (8.55,0)   {4. BASELINE};
\node[st=greenhl]   (e) at (11.4,0)   {5. FEATURES};
\node[st=purplehl]  (f) at (0,-2.55)  {6. MODEL};
\node[st=accent]    (g) at (2.85,-2.55){7. EVALUATE\\{\tiny \& threshold}};
\node[st=lensSEC]   (h) at (5.7,-2.55) {8. FAIRNESS\\{\tiny \& explain}};
\node[st=lensIOT]   (i) at (8.55,-2.55){9. DEPLOY\\{\tiny plan}};
\node[st=lensPM]    (j) at (11.4,-2.55){10. MONITOR\\{\tiny plan}};

\foreach \x/\y in {a/b,b/c,c/d,d/e}{\draw[arr] (\x) -- (\y);}
\draw[arr] (e.south) .. controls +(0,-1.1) and +(0,1.1) .. (f.north);
\foreach \x/\y in {f/g,g/h,h/i,i/j}{\draw[arr] (\x) -- (\y);}

\node[ch] at (0,-1.1)     {Ch.\ 2 --- decision, unit, target, timing, criterion, baseline};
\node[ch] at (2.85,-1.1)  {Ch.\ 3, 7 --- quality, missingness, leakage};
\node[ch] at (5.7,-1.1)   {Ch.\ 8 --- distributions, relationships, confounders};
\node[ch] at (8.55,-1.1)  {Ch.\ 11 --- beat it or stop};
\node[ch] at (11.4,-1.1)  {Ch.\ 15, 19 --- windows, lags, recency, trend};
\node[ch] at (0,-3.65)    {Ch.\ 12--18 --- simplest that works};
\node[ch] at (2.85,-3.65) {Ch.\ 14 --- PR curve, cost-based threshold};
\node[ch] at (5.7,-3.65)  {Ch.\ 28 --- subgroups, counterfactuals};
\node[ch] at (8.55,-3.65) {Ch.\ 24, 27 --- where, how, rollback};
\node[ch] at (11.4,-3.65) {Ch.\ 27 --- four layers, drift, retrain};
\end{tikzpicture}}%
{The ten-step capstone process, with the chapter that supplies each step. If you can carry
a problem through all ten and justify each choice, you have completed this
course.}{capstone-steps}

\section{The Four Deliverables}

\rowcolors{2}{white}{defbg}
\begin{longtable}{@{}L{3.4cm}L{4.4cm}L{6.2cm}@{}}
\toprule
\rowcolor{primary!15}
\textbf{Artefact} & \textbf{Audience} & \textbf{Must contain} \\
\midrule
\endhead
Technical report (8--12 pp) & Peer / examiner & Framing, audit, method, validation scheme, results with uncertainty, subgroup performance, limitations (\chref{research}) \\
Reproducible repository & Anyone, in a year & Pinned environment, fixed seeds, one command that regenerates every figure (\chref{research}) \\
One-page decision brief & The sponsor & Bottom line up front, cost, recommendation, the one caveat that changes the decision (\chref{communication}) \\
Model card & The record & Intended and out-of-scope use, data, subgroup metrics, ethical review, contest route (\chref{ethics}) \\
\bottomrule
\end{longtable}

\section{Marking Rubric --- Assess Yourself First}

\rowcolors{2}{white}{lightbg}
\begin{longtable}{@{}L{3.6cm}L{1.6cm}L{4.3cm}L{4.3cm}@{}}
\toprule
\rowcolor{primary!15}
\textbf{Criterion} & \textbf{Weight} & \textbf{Strong work} & \textbf{Weak work} \\
\midrule
\endhead
Framing & 15\% & All six questions answered; baseline identified before modelling & ``Predict X'' with no decision or baseline \\
Data handling & 15\% & Audited, leakage checked, missingness mechanism reasoned about & Loaded and modelled \\
Method choice & 15\% & Simplest adequate model, justified by a constraint & Most complex model available, unjustified \\
Evaluation & 20\% & Right metric for the base rate; threshold from costs; validation matches deployment & Accuracy on a random split \\
Ethics \& fairness & 15\% & Subgroup metrics, criterion chosen and defended, contest route & Mentioned in a closing paragraph \\
Operations & 10\% & Deployment tier justified; four-layer monitoring; retraining trigger & ``Then we deploy it'' \\
Communication & 10\% & Conclusion first; uncertainty in natural frequencies; three audiences served & Method narrated chronologically \\
\bottomrule
\end{longtable}

\begin{alertbox}[The Three Things That Sink Capstone Projects]
\begin{tightnum}
  \item \textbf{An unbelievably good result.} 99\% accuracy on a hard problem is
        leakage until proven otherwise (\chref{wrangling}). Examiners check this
        first, and finding it yourself is worth more marks than hiding it.
  \item \textbf{Accuracy on an imbalanced problem.} Reporting 97\% where 97\% is the
        majority class demonstrates that \chref{evaluation} was not understood.
  \item \textbf{No baseline.} Without it, no result can be shown to be worth
        anything, and the entire project is unassessable.
\end{tightnum}
\end{alertbox}

%---------------------------------------------------------------------
\begin{fourlenses}{The Same Project, Moved}
\lensLA{\emph{If you built here and moved elsewhere:} your interpretable model and
counterfactual explanations were driven by the fairness stakes of modelling people.
Moving to IoT, that constraint disappears and you may use a black box freely ---
which is a genuine gain in accuracy, obtained by leaving the hardest constraint
behind. State that honestly rather than claiming the model improved.}

\lensPM{\emph{If you built here:} your entire design was shaped by $n \approx 40$.
Moving to security, you gain billions of rows and lose the ability to inspect them,
so cross-validation stops being the bottleneck and label quality becomes it.
Conversely, moving \emph{into} project management from any other domain, expect to
discard your model class entirely --- ensembles cannot be justified on 40 rows.}

\lensIOT{\emph{If you built here:} you had abundant data and almost no labels, so
you used anomaly detection. Moving to learning analytics, labels appear but rows
collapse from millions to thousands, and --- decisively --- the entity becomes a
person with rights. Every architectural freedom you enjoyed at the edge is replaced
by an accountability requirement.}

\lensSEC{\emph{If you built here:} you designed against an adversary, chose features
by manipulation cost, and worked backwards from analyst capacity. Moving to IoT, the
adversary vanishes and you may use the cheap, obvious features you previously had to
avoid. Moving \emph{into} security from anywhere else, the reverse holds, and it is
the single largest adjustment in this book: you must now assume your model will be
studied and attacked.}
\end{fourlenses}

%---------------------------------------------------------------------
\begin{lab}[The Capstone Skeleton]
\begin{lstlisting}[language=Python]
"""capstone.py -- one command regenerates every figure in the report."""
import random, numpy as np
SEED = 42
random.seed(SEED); np.random.seed(SEED)

# ---- 1. FRAME (Ch 2) -- written down BEFORE any data is loaded --------
FRAME = {
    "decision":  "Which 20% of entities receive the limited intervention.",
    "unit":      "one entity per week",
    "target":    "bad outcome within the next 30 days",
    "timing":    "features from strictly before the prediction point",
    "criterion": "recall >= 0.70 at a flag rate <= 0.20",
    "baseline":  "current rule; measure it FIRST",
}

# ---- 2-3. AUDIT and EXPLORE (Ch 3, 7, 8) ------------------------------
audit(df)                     # completeness, validity, uniqueness, volume
assert not leaks(features, cutoff=PREDICTION_TIME)   # Ch 7: the fatal error

# ---- 4. BASELINE FIRST (Ch 11). Never skip. ---------------------------
base = current_rule(df)
print("baseline:", score(base))          # everything is measured against this

# ---- 5-6. FEATURES and MODEL (Ch 15, 19) ------------------------------
X = build_features(df)        # windows, lags, recency, trend, per-entity z-scores
model = fit(X, y, cv=correct_cv_for_domain())   # walk-forward / group / stratified

# ---- 7. EVALUATE: right metric, cost-based threshold (Ch 14) ----------
report_pr_curve(model, X_te, y_te)                   # NOT roc if positives are rare
thr = threshold_from_costs(cost_fp=COST_FP, cost_fn=COST_FN)

# ---- 8. FAIRNESS (Ch 28) ---------------------------------------------
fairness_report(y_te, y_pred, y_score, group=protected_attribute)
assert recall_gap() < 0.05, "subgroup recall gap too large -- do not deploy"

# ---- 9-10. DEPLOY and MONITOR plans (Ch 24, 27) ----------------------
write_model_card(FRAME, metrics, subgroup_metrics, limitations, contest_route)
write_monitoring_plan(layers=["inputs", "predictions", "outcomes", "business"])
\end{lstlisting}

\textbf{Notice what this skeleton enforces:} the frame exists before the data, the
baseline exists before the model, the leakage check is an assertion rather than an
intention, and the fairness gap can block deployment. Those four lines are the
difference between a student project and a professional one.
\end{lab}

%---------------------------------------------------------------------
\begin{checkpoint}
\begin{tightnum}
  \item Your capstone model reaches 0.97 accuracy on a problem where 6\% of cases are
        positive. What are the first two things you check, in order?
  \item You built for IoT and are asked to move the approach to learning analytics.
        Name the three constraints that change and what each forces you to do
        differently.
  \item Which single deliverable of the four would a sponsor read, and what must be
        in its first two sentences?
\end{tightnum}
\end{checkpoint}

%---------------------------------------------------------------------
\begin{chaptersummary}
\begin{tight}
  \item All four course domains are the same problem: early detection of rare bad
        outcomes from behavioural traces under a capacity constraint.
  \item What differs between them is the base rate, label availability,
        ground-truth lag, the presence of an adversary and the human stakes --- not
        the algorithm.
  \item The ten-step process runs framing $\rightarrow$ audit $\rightarrow$ explore
        $\rightarrow$ baseline $\rightarrow$ features $\rightarrow$ model
        $\rightarrow$ evaluate $\rightarrow$ fairness $\rightarrow$ deploy
        $\rightarrow$ monitor.
  \item Deliver four artefacts: technical report, reproducible repository, one-page
        decision brief, model card.
  \item Projects fail on leakage, on accuracy under imbalance, and on the absence of
        a baseline --- all three preventable in the first week.
  \item The mark of competence is not building the most complex model, but choosing
        the simplest one the constraints permit and saying which constraint forced
        each choice.
\end{tight}
\end{chaptersummary}

\begin{reviewq}
  \item \textbf{(Short)} State, in one sentence each, the four properties that
        distinguish the course's domains from one another.
  \item \textbf{(Short)} Why is ``we used gradient boosting because it was most
        accurate'' a weaker justification than ``we used logistic regression because
        advisers must be able to explain the decision to a student''?
  \item \textbf{(Short)} Give the three checks an examiner performs first on a
        suspiciously good result.
  \item \textbf{(Applied)} Complete the full ten-step process for one domain of your
        choice. Produce all four deliverables. Then write the one-page transfer
        analysis required by the brief.
  \item \textbf{(Applied)} Take a project you have already completed --- in this or
        another module --- and mark it against the rubric above. Identify the two
        criteria on which it scored worst and write what you would do differently.
  \item \textbf{(Applied)} You are given 40 sprints of project data and asked for a
        completion-risk model. Write the honest one-paragraph response you would send
        the sponsor, including what you can and cannot deliver, and what you would do
        instead.
\end{reviewq}

\begin{center}
\vspace{6pt}
\begin{tcolorbox}[enhanced, width=0.94\textwidth, colback=primary!5,
  colframe=primary, boxrule=0pt, leftrule=5pt, arc=3pt, top=8pt, bottom=8pt]
\small
\textbf{\color{primary}A closing word.} You began this handout being told that data
science turns data into understanding and action. Thirty-two chapters later, the
part that turned out to be hard was never the modelling. It was framing the question
so that an answer would matter, getting data honest enough to trust, choosing a
metric that reflected the real cost of being wrong, and explaining the result to
someone who had to act on it. The algorithms are the easy part, and they are
freely available to everyone. The judgement is not.
\end{tcolorbox}
\end{center}
